\authorchapter{Reflexion}{Dražen Plavšić, Ereny Micail}
\label{ch:reflexion}

\authorsection{Gemeinsame Reflexion}{Dražen Plavšić, Ereny Micail}
Die im Zeitplan (Abschnitt~\ref{ch:projektmanagement}) definierten Meilensteine wurden größtenteils eingehalten; einzelne Phasen verzögerten sich geringfügig, insbesondere durch die nachträgliche Integration der Usability-Test-Ergebnisse gegen Projektende, ohne dass dies den finalen Abgabetermin gefährdete.
%------------------------------------
% DRAŽEN
%------------------------------------
Die Zusammenarbeit im Team verlief insgesamt problemlos. Auftretende Schwierigkeiten konnten gemeinsam und ohne größere Abstimmungen gelöst werden. Beide Seiten haben sich bei Bedarf gegenseitig unterstützt und Aufgaben flexibel übernommen. Auch nicht klar zugewiesene Themenbereiche wurden fair verteilt, ohne dass es dabei zu größeren Diskussionen oder Konflikten kam.

Ein weiteres Verbesserungspotenzial zeigte sich in der zeitlichen Platzierung der Tests. Durch die anfängliche Orientierung am Wasserfallmodell wurden automatisierte Tests erst in einer späteren Projektphase entwickelt. Dadurch mussten Fehler häufig erst nach der Implementierung korrigiert werden, was zu zusätzlichen Bugfix-Schleifen führte. Ein früherer Einsatz von Test-Driven Development hätte in diesem Zusammenhang vermutlich dazu beigetragen, Probleme früher zu erkennen und den Korrekturaufwand zu reduzieren.

\authorsection{Evaluierung alternativer Anonymisierungsstrategien}{Dražen Plavšić}
\label{sec:live_chat_alternativen}
Während die Zuweisung auf einen zentralen Ghost-User (siehe \autoref{sec:dsgvo}) eine praktische und im Sinne der Terminologie der \gls{ENISA} irreversible Anonymisierungslösung darstellt \cite[S.~9--10]{enisa2019pseudonymisation}, wurde im Rahmen der Architekturplanung auch die sogenannte \textit{In-Place-Anonymisierung} in Betracht gezogen.

Bei diesem Ansatz wird der eigentliche Datenbankeintrag des ausscheidenden Nutzers nicht physisch gelöscht (Hard Delete). Der Datensatz, und damit der originale Primärschlüssel (\texttt{user\_id}), bleibt erhalten. Stattdessen werden sämtliche personenbezogenen Felder (wie Vorname, Nachname oder Avatar) durch generische Platzhalter oder \texttt{NULL}-Werte überschrieben.

Um Datenbank-Restriktionen wie \texttt{UNIQUE}-Constraints (z.\,B. bei der E-Mail-Adresse) nicht zu verletzen, werden diese identifizierenden Felder durch einen kryptografischen Hash (beispielsweise SHA-256) ersetzt. Hierbei ist jedoch kritisch anzumerken, dass der Einsatz einfacher Hash-Funktionen laut der \gls{ENISA} für Dictionary Attacks anfällig ist. Dementsprechend gilt ein ungesalzenes Hashing (\textit{Plain Hashing}) datenschutzrechtlich lediglich als Pseudonymisierung und explizit nicht als Anonymisierung \cite[S.~22, 33]{enisa2019pseudonymisation}.

Dieser Ansatz bietet dennoch einen signifikanten Vorteil. In der Nachrichtenhistorie bleibt eindeutig nachvollziehbar, ob aufeinanderfolgende Nachrichten von derselben gelöschten Person oder von unterschiedlichen ehemaligen Nutzern stammen (z.\,B. \glqq{}Ehemaliger Nutzer 1\grqq{} vs. \glqq{}Ehemaliger Nutzer 2\grqq{}). Bei einem zentralen Ghost-User hingegen verschmelzen alle gelöschten Identitäten zu einer einzigen, was die Nachvollziehbarkeit der Kommunikationshistorie dauerhaft einschränkt.

Obwohl die In-Place-Methode die Kollaborationshistorie präziser abbildet, muss berücksichtigt werden, dass das Behalten des originalen Primärschlüssels die Gefahr einer Reidentifizierung (durch Zusatzinformationen) birgt und somit den Vorgaben von Art. 17 DSGVO widersprechen könnte. Da es laut \gls{ENISA} keine universelle Standardlösung für derartige Datenbankkonflikte gibt, müssen Datenschutzmaßnahmen stets auf einem risikobasierten Ansatz beruhen\cite[S.~42]{enisa2019pseudonymisation}.

\authorsection{Reflexion der Backend-Architektur: Table-Module vs. Domain-Model}{Dražen Plavšić}
\label{sec:table_modul_pattern_vs_domain_model}
Im Rückblick auf die Entwicklung des Backends lässt sich feststellen, dass der gewählte Ansatz des \textit{Table-Module-Pattern} (siehe \autoref{sec:backend_architektur}) zwar eine klare und nachvollziehbare Grundstruktur bietet, jedoch in der praktischen Umsetzung systemweit zu architektonischen Schwächen geführt hat. Martin Fowler beschreibt das \textit{Table-Module} als ein Muster, das sich für einfache \gls{CRUD}-Anwendungen eignet, jedoch bei steigender Komplexität der Geschäftslogik an seine Grenzen stößt\cite[S.~125--131]{fowler2002patterns}.

Dieses grundlegende Problem betrifft prinzipiell die gesamte Applikation, zeigt sich jedoch besonders deutlich bei der Taskverwaltung. Operationen wie das Erstellen, Aktualisieren oder Löschen eines Tasks erfordern das gleichzeitige Befüllen der \texttt{task\_activity\_log}-Tabelle, um einen konsistenten Systemzustand zu gewährleisten. Da diese Operationen systembedingt auf separate Module und Controller verteilt sind, werden sie nicht zwingend als atomare Datenbanktransaktion ausgeführt. Im Fehlerfall, etwa wenn der Task erfolgreich aktualisiert wird, der zugehörige Activity-Log-Eintrag jedoch fehlschlägt, kann dies zu inkonsistenten Daten in der Datenbank führen.

Ein auf \gls{DDD} basierender Architekturansatz hätte hier für das gesamte Backend eine bessere Alternative dargestellt \cite[S.~68--77]{evans2004domain}. Dabei würde die Modulstruktur nicht nach Datenbanktabellen, sondern nach fachlichen Domänen und Geschäfts\-prozessen organisiert. Bezogen auf das vorherige Beispiel wäre ein einziger \texttt{TaskController} in diesem Fall verantwortlich für alle taskbezogenen Operationen, inklusive der Pflege des Activity-Logs innerhalb einer gemeinsamen Transaktion. Dies hätte die Datenkonsistenz über alle Modulgrenzen hinweg strukturell sichergestellt und die Verantwortlichkeiten systemweit klarer nach dem \textit{Single-Responsibility-Prinzip} abgegrenzt \cite[S.~56--61]{martin2017clean}. In DDD-Terminologie würde die Taskverwaltung dabei einen eigenen \gls{bounded-context} bilden, innerhalb dessen der \texttt{TaskController} als \gls{aggregate-root} sämtliche Konsistenzregeln kapselt.

\authorsection{Persönliche Reflexion}{Ereny Micail}
Im Rückblick auf das Projekt lag die größte persönliche Herausforderung weniger in einzelnen fachlichen Details als im Zeitmanagement und der Priorisierung paralleler Arbeitspakete. Da mehrere thematisch unterschiedliche Module -- Authentifizierung, Rollen- und Berechtigungslogik, Kommunikation, Benachrichtigungen und Kalender (siehe Abschnitte~\ref{ch:implementierung_kernfunktionen}) -- in überlappenden Zeiträumen bearbeitet wurden, war es phasenweise schwierig, den jeweils dringendsten Punkt korrekt einzuschätzen. Besonders deutlich zeigte sich dies gegen Projektende, als die im Usability-Test (Abschnitt~\ref{sec:usability_test}) identifizierten Probleme kurzfristig in die laufende Planung integriert werden mussten, ohne dass dafür ursprünglich ausreichend Pufferzeit eingeplant war.

Die wichtigste persönliche Erkenntnis aus dem Projekt ist, dass uns zu Beginn die notwendige Erfahrung fehlte, um den Zeitaufwand für die einzelnen Aufgaben realistisch einzuschätzen. Dies führte dazu, dass viele Arbeitspakete deutlich mehr Zeit in Anspruch nahmen als ursprünglich kalkuliert. Erschwerend kam hinzu, dass private Verpflichtungen die Einhaltung des ohnehin knappen Zeitplans zusätzlich erschwerten. Rückblickend wäre es weitaus produktiver gewesen, wenn wir uns bereits vor der Planungsphase mit spezifischen Projektmanagement-Methoden zur Zeiteinteilung vertraut gemacht hätten, um die Struktur und Dauer der Arbeitsschritte fundierter festlegen zu können.

Dieser Mangel an Vorkenntnissen beschränkte sich jedoch nicht nur auf die Zeitplanung: Auch der Einsatz automatisierter Tests wurde dadurch in Mitleidenschaft gezogen – uns fehlte schlicht das nötige Wissen, um deren Bedeutung frühzeitig zu erkennen und sie von Beginn an konsequent einzusetzen. Wie unsere gemeinsame Reflexion zeigt, hätte ein früherer Einsatz den späteren Korrekturaufwand erheblich verringert – in beiden Fällen wäre also ein Wechsel von spontanem, reagierendem Arbeiten hin zu vorausschauendem Arbeiten entscheidend gewesen.

Für zukünftige Projekte nehme ich mir deshalb konkret vor, Aufgaben von Beginn an in kleinere, klar abgegrenzte Teilschritte zu gliedern, um den tatsächlichen Aufwand realistischer einschätzen zu können. Zeitpuffer werde ich nicht erst als Reaktion auf eingetretene Verzögerungen einplanen, sondern bereits in der initialen Planung als festen Bestandteil verankern. Darüber hinaus möchte ich mich vor dem Start eines Projekts gezielt mit relevanten Projektmanagement-Methoden auseinandersetzen, um strukturierter und vorausschauender vorgehen zu können – und nicht erneut in ein reaktives Arbeiten zu verfallen, das unnötigen Druck erzeugt.